\tikzexternaldisable
\begin{center}
\begin{tikzpicture}[scale=1.1, >=stealth]

  % Grilla de fondo sutil
  \draw[very thin, gray!20] (-0.5,-0.5) grid (4.5,3.5);

  % Ejes Cartesianos
  \draw[->, thick] (-0.5, 0) -- (4.8, 0) node[right] {$x$};
  \draw[->, thick] (0, -0.5) -- (0, 3.8) node[above] {$y$};

  % Marcas en los ejes
  \foreach \x in {1,2,3,4}
    \draw (\x, 0.08) -- (\x, -0.08) node[below, font=\small] {$\x$};
  \foreach \y in {1,2,3}
    \draw (0.08, \y) -- (-0.08, \y) node[left, font=\small] {$\y$};

  % Relleno de la región encerrada por la curva C
  \fill[blue!15, opacity=0.7] (0,0) -- (2,3) -- (4,2) -- cycle;

  % Trayectoria poligonal C con orientación
  % Segmento AB (C1)
  \draw[line width=1.1pt, blue!70!black, postaction={decorate, decoration={
    markings, mark=at position 0.55 with {\arrow{>}}}}] (0,0) -- (2,3) 
    node[midway, above left, font=\small] {$C_1$};

  % Segmento BC (C2) -> Se cambió 'above right' por 'above, sloped' o desplazamiento fino
  \draw[line width=1.1pt, blue!70!black, postaction={decorate, decoration={
    markings, mark=at position 0.55 with {\arrow{>}}}}] (2,3) -- (4,2) 
    node[midway, above right=1pt, font=\small] {$C_2$};

  % Segmento CA (C3)
  \draw[line width=1.1pt, blue!70!black, postaction={decorate, decoration={
    markings, mark=at position 0.55 with {\arrow{>}}}}] (4,2) -- (0,0) 
    node[midway, below, font=\small] {$C_3$};

  % Puntos/Vértices A, B, C
  \fill[blue!80!black] (0,0) circle (2pt) node[below left, font=\small\bfseries] {$A(0,0)$};
  \fill[blue!80!black] (2,3) circle (2pt) node[above, font=\small\bfseries] {$B(2,3)$};
  \fill[blue!80!black] (4,2) circle (2pt) node[right, font=\small\bfseries] {$C(4,2)$};

  % Etiquetas D y C
  \node[color=blue!80!black, font=\large] at (2, 2) {$D$};
  % La etiqueta C de la frontera total se coloca cerca del tramo C1 o arriba
  \node[color=blue!70!black, font=\large] at (0.9, 2.5) {$C$};

\end{tikzpicture}
\end{center}
\tikzexternalenable